\documentclass[11pt,a4paper]{article}
\usepackage[margin=2.5cm]{geometry}
\usepackage{amsmath,amssymb}
\usepackage{booktabs}
\usepackage{hyperref}
\usepackage{listings}
\lstset{basicstyle=\ttfamily\small,columns=fullflexible,frame=single,language=Fortran}

\newcommand{\bx}{\mathbf{x}}
\newcommand{\bv}{\mathbf{v}}
\newcommand{\ba}{\mathbf{a}}
\newcommand{\bj}{\mathbf{j}}
\newcommand{\T}{\mathsf{T}}
\newcommand{\bO}{\boldsymbol{\Omega}}
\newcommand{\kmskpc}{\mathrm{km\,s^{-1}\,kpc^{-1}}}

\title{The Galactic tidal field in HENDOS\\[4pt]
\large HErmite $N$-body Dynamics Of Systems}
\author{Sheffield star formation and dynamics group}
\date{\today}

\begin{document}
\maketitle

\section{Overview}

The cluster is assumed to move on a circular orbit in an axisymmetric
Galaxy. The integration is carried out in a frame centred on the
cluster's guiding centre and rotating with the orbit, and the Galactic
potential is expanded to second order about the guiding centre (Hill's
approximation; e.g.\ Chandrasekhar 1942; Heggie \& Hut 2003; Renaud,
Gieles \& Boily 2011). The external field then enters the equations of
motion as a constant, symmetric \emph{tidal tensor} $\T$ plus the
Coriolis force. Because $\T$ is constant in the rotating frame, the
Jacobi energy is conserved and can be used to monitor the accuracy of
the integration in the same way as the energy of an isolated cluster.

Three options are available through the parameter \texttt{tidal} in the
\texttt{\&tides} section of the parameter file:
\begin{itemize}
\item \texttt{'none'}: an isolated cluster ($\T=0$, $\Omega=0$);
\item \texttt{'solar'}: the solar-neighbourhood field, built from the
  Oort constants and the local mass density (Section~\ref{sec:solar});
\item \texttt{'tensor'}: any user-supplied tensor $\T$ and frame
  rotation rate $\Omega$.
\end{itemize}

\section{Coordinate frame}

The frame rotates with angular velocity $\bO = \Omega\,\hat{\mathbf z}$
about the cluster centre, with
\begin{itemize}
\item $x$ pointing away from the Galactic centre,
\item $y$ pointing in the direction of Galactic rotation,
\item $z$ completing a right-handed set, i.e.\ along the Galaxy's spin
  (towards the south Galactic pole).
\end{itemize}
Initial positions and velocities are read, and snapshots are written,
in this frame. Velocities are therefore velocities \emph{relative to
the rotating frame}. A cluster that is non-rotating in an inertial frame
appears to rotate at $-\Omega$ in this frame, but for the solar value
($\Omega \simeq 27\,\kmskpc$) this amounts to only
$\sim 0.03\,\mathrm{km\,s^{-1}}$ per pc of radius, which is negligible
compared with typical velocity dispersions.

\section{Equations of motion}

For star $i$ at position $\bx_i$ with velocity $\bv_i$ in the rotating
frame,
\begin{equation}
  \ddot{\bx}_i = \ba_{\mathrm{grav},i} + \T\bx_i - 2\bO\times\bv_i ,
  \label{eq:eom}
\end{equation}
where $\ba_{\mathrm{grav},i}$ is the gravitational acceleration from
the other stars. The term $\T\bx$ is the tidal acceleration including
the centrifugal term, and $-2\bO\times\bv$ is the Coriolis
acceleration. In components, with $\bO=\Omega\hat{\mathbf z}$,
\begin{equation}
  -2\bO\times\bv = 2\Omega\,(v_y,\,-v_x,\,0).
\end{equation}

The effective tensor is
\begin{equation}
  \T = \T_{\mathrm{G}} + \Omega^2\,\mathrm{diag}(1,1,0),
  \qquad
  (\T_{\mathrm{G}})_{kl} = -\left.\frac{\partial^2\Phi_{\mathrm{G}}}
  {\partial x_k\,\partial x_l}\right|_{\mathrm{gc}},
\end{equation}
where $\Phi_{\mathrm G}$ is the Galactic potential and the derivatives
are taken at the guiding centre. The second term is the centrifugal
contribution of the rotating frame. $\T$ must be symmetric; if a
non-symmetric tensor is supplied, the code warns and uses its
symmetric part $(\T+\T^{\mathsf T})/2$.

\section{The solar-neighbourhood field}
\label{sec:solar}

For an axisymmetric potential with circular speed $v_c(R)$ the angular
velocity of the frame and the derivatives of $\Phi_{\mathrm G}$ at the
guiding centre can be written in terms of the Oort constants
$A = -\tfrac12 R\,\mathrm{d}\Omega/\mathrm{d}R$ and $B = A-\Omega$:
\begin{align}
  \Omega &= A - B, \\
  \frac{\partial^2\Phi_{\mathrm G}}{\partial R^2} &= \Omega^2 - 4A\Omega, \\
  \frac{1}{R}\frac{\partial\Phi_{\mathrm G}}{\partial R} &= \Omega^2, \\
  \frac{\partial^2\Phi_{\mathrm G}}{\partial z^2} &\equiv \nu^2
     = 4\pi G\rho_0 + 2\left(A^2 - B^2\right),
\end{align}
where the last line follows from Poisson's equation in cylindrical
coordinates and $\rho_0$ is the total (stars, gas and dark matter)
mass density at the guiding centre. Adding the centrifugal term gives
\begin{equation}
  \T = \mathrm{diag}\left(4A(A-B),\; 0,\; -\left[4\pi G\rho_0 +
  2\left(A^2-B^2\right)\right]\right).
\end{equation}
The equations of motion then reduce to the familiar Hill equations
\begin{align}
  \ddot x &= a_{\mathrm{grav},x} + 4A\Omega\, x + 2\Omega\,\dot y, \\
  \ddot y &= a_{\mathrm{grav},y} - 2\Omega\,\dot x, \\
  \ddot z &= a_{\mathrm{grav},z} - \nu^2 z .
\end{align}

The default values are given in Table~\ref{tab:solar}. They can be
changed in the parameter file.

\begin{table}[h]
\centering
\begin{tabular}{lll}
\toprule
Quantity & Value & Source \\
\midrule
$A$ & $15.3\,\kmskpc$ & Bovy (2017) \\
$B$ & $-11.9\,\kmskpc$ & Bovy (2017) \\
$\rho_0$ & $0.097\,\mathrm{M_\odot\,pc^{-3}}$ & McKee, Parravano \& Hollenbach (2015) \\
\midrule
$\Omega = A-B$ & $27.2\,\kmskpc$ & \\
$T_{xx} = 4A\Omega$ & $1664.6\,(\kmskpc)^2$ & \\
$T_{zz} = -\nu^2$ & $-5427.5\,(\kmskpc)^2$ & \\
\bottomrule
\end{tabular}
\caption{Default solar-neighbourhood parameters and the resulting
tidal field.}
\label{tab:solar}
\end{table}

\subsection{Jacobi radius}
For a cluster of mass $M$ the Jacobi (tidal) radius along the $x$-axis
is
\begin{equation}
  r_{\mathrm J} = \left(\frac{GM}{T_{xx}}\right)^{1/3}
  = \left(\frac{GM}{4A\Omega}\right)^{1/3},
\end{equation}
which the code prints at the start of any run with $T_{xx}>0$. For
$M = 1000\,\mathrm{M_\odot}$ in the solar neighbourhood
$r_{\mathrm J} \simeq 13.7$\,pc.

\section{Energy}

The Coriolis force does no work and $\T$ is constant and symmetric, so
the tidal acceleration derives from the potential
\begin{equation}
  \Phi_{\mathrm{tid}}(\bx) = -\tfrac12\,\bx^{\mathsf T}\,\T\,\bx .
\end{equation}
The conserved quantity is the Jacobi energy
\begin{equation}
  E_{\mathrm J} = \sum_i \tfrac12 m_i |\bv_i|^2
   \;-\; \sum_{i<j}\frac{G m_i m_j}{\left(r_{ij}^2+\epsilon^2\right)^{1/2}}
   \;+\; \sum_i m_i\,\Phi_{\mathrm{tid}}(\bx_i),
\end{equation}
with velocities measured in the rotating frame. When a tide is
switched on, the energy and relative energy error printed by the code
are those of $E_{\mathrm J}$. The virial ratio printed is still
$-K/W$ using the gravitational potential energy $W$ alone.

\section{Implementation in the Hermite scheme}

\subsection{Acceleration and jerk}
The fourth-order Hermite scheme needs both the acceleration and its
time derivative (the jerk) for each star. Differentiating
equation~(\ref{eq:eom}) gives the tidal contribution to the jerk,
\begin{equation}
  \bj_{\mathrm{tid}} = \T\bv - 2\bO\times\ba ,
\end{equation}
where $\ba$ is the \emph{total} acceleration (gravitational plus tidal
plus Coriolis). In the force routine the gravitational sum over all
other stars is done first; the tidal and Coriolis accelerations are
then added; and only then are the tidal jerk terms added, so that they
use the complete acceleration:
\begin{lstlisting}
  ax = ax + T11*x + T12*y + T13*z + 2*omega*vy
  ay = ay + T21*x + T22*y + T23*z - 2*omega*vx
  az = az + T31*x + T32*y + T33*z
  jx = jx + T11*vx + T12*vy + T13*vz + 2*omega*ay
  jy = jy + T21*vx + T22*vy + T23*vz - 2*omega*ax
  jz = jz + T31*vx + T32*vy + T33*vz
\end{lstlisting}
The cost is a few operations per active star per step, negligible
compared with the $O(N)$ gravitational sum.

\subsection{Predictor order}
In the standard Hermite scheme positions are predicted to third order
and velocities to second order,
\begin{align}
  \bx_{\mathrm p} &= \bx_0 + \bv_0\Delta t + \tfrac12\ba_0\Delta t^2
    + \tfrac16\bj_0\Delta t^3, \\
  \bv_{\mathrm p} &= \bv_0 + \ba_0\Delta t + \tfrac12\bj_0\Delta t^2 .
\end{align}
For purely gravitational forces, which depend only on positions, this
is sufficient for fourth-order accuracy. The Coriolis force, however,
depends on velocity, and an $O(\Delta t^3)$ error in $\bv_{\mathrm p}$
then gives an $O(\Delta t^3)$ error in the force evaluated at the
predicted state. This reduces the whole scheme to third order. Tests
confirmed this: with a strong tide the energy error fell by a factor of
$\approx 8$, not 16, for each halving of the timestep.

The code therefore raises the order of the prediction used for the
force evaluation by adding the second and third derivatives of the
acceleration, $\ba^{(2)}$ and $\ba^{(3)}$, from the end of each star's
previous step:
\begin{align}
  \bx_{\mathrm p} &\leftarrow \bx_{\mathrm p}
     + \tfrac1{24}\ba^{(2)}\Delta t^4 + \tfrac1{120}\ba^{(3)}\Delta t^5, \\
  \bv_{\mathrm p} &\leftarrow \bv_{\mathrm p}
     + \tfrac16\ba^{(2)}\Delta t^3 + \tfrac1{24}\ba^{(3)}\Delta t^4 .
\end{align}
Here $\ba^{(2)} = \mathbf{f}_2 + \mathbf{f}_3\Delta t_{\mathrm{prev}}$
and $\ba^{(3)} = \mathbf{f}_3$ come from the previous corrector step
(arrays \texttt{f2e}, \texttt{f3e}); they are zero on the first step.
The corrector itself is unchanged and still starts from the
third-order prediction, because it adds the $\mathbf f_2$ and
$\mathbf f_3$ terms itself; using the higher-order prediction there
would count those terms twice. The improved prediction is also used
for the non-active stars, which helps the accuracy of block timesteps
in general.

\subsection{Timesteps}
No special timestep criterion is needed for the tide. The Aarseth
criterion uses the total acceleration and its derivatives, so stars far
from the cluster, whose motion is dominated by the tide, automatically
take steps set by the tidal frequencies. The largest step any star can
take is the output interval \texttt{tout}.

\subsection{Units}
Internally the code uses $G=1$, $\mathrm{M_\odot}$ and pc, giving a
velocity unit of $0.06558\,\mathrm{km\,s^{-1}}$ and a time unit of
$14.91$\,Myr. The tidal parameters are read in $\kmskpc$ (rates) and
$(\kmskpc)^2$ (tensor components), and converted with
\begin{equation}
  1\,\kmskpc = \frac{1}{1000\,v_{\mathrm{unit}}}
  \ \text{internal units} ,
\end{equation}
where $v_{\mathrm{unit}}$ is in $\mathrm{km\,s^{-1}}$. Because the
internal mass and length units are $\mathrm{M_\odot}$ and pc with
$G=1$, the density term $4\pi G\rho_0$ is simply $4\pi\rho_0$ with
$\rho_0$ in $\mathrm{M_\odot\,pc^{-3}}$. The user never sees internal
units: the tensor and $\Omega$ are echoed in $\kmskpc$ at the start of
each run.

\section{Tests}

The scheme was checked with a line-by-line Python version of the code,
using a 300-star Plummer sphere of $1000\,\mathrm{M_\odot}$ with a
Plummer radius of 1\,pc, integrated for 8\,Myr with $\eta = 0.02$:
\begin{itemize}
\item The analytic tidal jerk agrees with a finite-difference
  derivative of the acceleration along the motion, with an error that
  falls linearly with the difference step, as it should.
\item With the solar tide, $|\Delta E_{\mathrm J}/E_{\mathrm J}| =
  1.7\times10^{-7}$, comparable to the isolated cluster
  ($5.6\times10^{-7}$ with the standard predictor).
\item With a tide 1000 times stronger than the solar value
  ($r_{\mathrm J}\simeq1.4$\,pc, so the cluster is strongly stripped),
  $|\Delta E_{\mathrm J}/E_{\mathrm J}| = 3.6\times10^{-6}$ with the
  higher-order predictor, compared with $2.5\times10^{-3}$ with the
  standard predictor, at the same cost.
\end{itemize}

\section{Parameter file}

\begin{lstlisting}
&tides
  tidal = 'solar'     ! 'none', 'solar' or 'tensor'
  oorta = 15.3        ! Oort A (km/s/kpc)
  oortb = -11.9       ! Oort B (km/s/kpc)
  rho0  = 0.097       ! local total density (Msun/pc^3)
  ! used only when tidal = 'tensor', units (km/s/kpc)^2 and km/s/kpc
  tij(1,1) = 1664.6,  tij(3,3) = -5427.5
  omegat   = 27.2
/
\end{lstlisting}

\section*{References}
\begin{itemize}\setlength\itemsep{0pt}
\item Bovy J., 2017, MNRAS, 468, L63
\item Chandrasekhar S., 1942, \emph{Principles of Stellar Dynamics},
  University of Chicago Press
\item Heggie D., Hut P., 2003, \emph{The Gravitational Million-Body
  Problem}, Cambridge University Press
\item Makino J., Aarseth S.~J., 1992, PASJ, 44, 141
\item McKee C.~F., Parravano A., Hollenbach D.~J., 2015, ApJ, 814, 13
\item Renaud F., Gieles M., Boily C.~M., 2011, MNRAS, 418, 759
\end{itemize}

\end{document}
